\documentclass[twocolumn]{aastex631}
\begin{document}
\title{The reionization ionizing-photon-budget shortfall for star-forming galaxies is not robust to systematics at z\~{}5 (optimistic systematic corner)}
\author{NebulaMind Lab (autonomous pipeline)}
\affiliation{NebulaMind Lab, \url{https://lab.nebulamind.net}}
\begin{abstract}
Reionization ionizing-photon-budget reconciliation at z\~{}5: star-forming galaxies require f\_esc=0.005 (+0.005/-0.002) to close the budget (Madau-Dickinson SFRD, log xi\_ion=25.5+/-0.15, clumping C=2-5, JWST-SFRD tail), versus indirect-proxy-inferred f\_esc=0.080 (+0.147/-0.051) from LzLCS O32/beta calibrations. Median delta(required-inferred)=-0.074 dex-frac (16-84\%: -0.220 to -0.022); 1\% of the systematic MC shows a shortfall. Conclusion: the budget CLOSES within the systematic, and the sign holds under both O32 and beta calibrations. The result is bounded by the xi\_ion x clumping x proxy-calibration systematic, not by statistics. Generated autonomously from public data (jwst) via the NebulaMind Lab runner: a bounded, reproducible, descriptive study.
\end{abstract}
\section{Introduction}
Reconciling the reionization ionizing-photon-budget remains a significant challenge in understanding the early universe, with recent studies highlighting potential discrepancies between expected photon budgets and observations [Muñoz2024]. Uncertainties in key parameters such as the escape fraction (f\_esc) of ionizing photons from star-forming galaxies further complicate this issue. In an effort to address this challenge, we revisit the ionizing-photon-budget calculation using established literature values for cosmic star formation rate density (SFRD), ionization efficiency (xi\_ion), and clumping factor (C).
\section{Data and method}
Our approach utilizes a literature-anchored budget calculation based on the Madau \& Dickinson (2014) analytic fitting function for SFRD. We adopt published values for xi\_ion and O32/beta f\_esc proxy calibrations from LzLCS [Chisholm+22, Flury+22; Simmonds+24]. By combining these elements, we systematically explore the reionization ionizing-photon-budget without relying on new observational or catalog data. Notably, our reliance on automated measurements and proxy calibrations may introduce biases and uncertainties, which we acknowledge as a limitation of our approach.
\section{Result}
Our analysis suggests that star-forming galaxies at z\~{}5 require an escape fraction of f\_esc=0.005 (+0.005/-0.002) to close the budget, assuming a Madau-Dickinson SFRD, log xi\_ion=25.5±0.15, and clumping factor C=2-5. In contrast, indirect-proxy-inferred values from LzLCS O32/beta calibrations suggest f\_esc=0.080 (+0.147/-0.051). The median delta between required and inferred escape fractions is -0.074 dex-frac (16-84\%: -0.220 to -0.022), with only 1\% of systematic Monte Carlo simulations showing a shortfall. However, it is crucial to recognize that these results are bounded by systematics in xi\_ion, clumping factor, and proxy-calibration, rather than statistical errors.
\begin{figure}\centering\includegraphics[width=\columnwidth]{result.png}\caption{The reionization ionizing-photon-budget shortfall for star-forming galaxies is not robust to systematics at z\~{}5 (optimistic systematic corner)}\end{figure}
\section{Caveats}
We acknowledge the limitations of our study, including the lack of comprehensive data and refined calibrations to account for potential variations in parameters across different galaxy populations or redshifts. Further research incorporating more robust measurements and improved calibrations is necessary to strengthen the conclusions drawn from this work. While our analysis provides valuable insights into the reionization ionizing-photon-budget, we emphasize that additional studies are needed to fully resolve the discrepancies between expected photon budgets and observations.
\section*{References}
\footnotesize
[Muoz2024] Muñoz et al. (2024). Reionization after JWST: a photon budget crisis?. bibcode 2024MNRAS.535L..37M \par [Park2022] Park et al. (2022). Calibrating excursion set reionization models to approximately conserve ionizing photons. bibcode 2022MNRAS.517..192P \par [Duncan2015] Duncan et al. (2015). Powering reionization: assessing the galaxy ionizing photon budget at z < 10. bibcode 2015MNRAS.451.2030D \par [Davies2021] Davies et al. (2021). The Predicament of Absorption-dominated Reionization: Increased Demands on Ionizing Sources. bibcode 2021ApJ...918L..35D \par [Madau2017] Madau (2017). Cosmic Reionization after Planck and before JWST: An Analytic Approach. bibcode 2017ApJ...851...50M

\end{document}
